Efficient information seeking requires questions that split the possibility space in half---the optimal strategy from binary search.
We investigate whether modern large language models can, in a zero-shot setting, generate yes/no questions that achieve this maximum-entropy split over a given set of items.
We evaluate \gptfull and \gptmini across 8 datasets (1,049 trials) with three prompt strategies: basic, explicit split instruction, and chain-of-thought (CoT).
We find that the capability exists but is not the default behavior: without explicit guidance, \gptfull achieves near-zero binary entropy on homogeneous sets (0.027 on fruits), yet with CoT prompting jumps to 0.977 on the same sets and 1.000 on structured category sets.
Surprisingly, \gptmini produces more naturally discriminative questions than \gptfull under basic prompting, though \gptfull surpasses it with explicit instructions.
Set composition is the primary difficulty factor: mixed-category sets yield 100\% perfect splits, while homogeneous sets (all fruits or animals) top out at 65--70\%.
These results demonstrate that modern LLMs possess strong implicit knowledge of item properties sufficient for optimal set partitioning, but this reasoning must be explicitly elicited through prompt design.
